% The passages replaced at v3.100 (AUDIT_v3.98 §13, applied 2026-09-25), verbatim, by tag.
% Not input by any master; kept for the record. New text: chapters/05_setup.tex, 06_results.tex.

% ===== [A1a platforms prose] chapters/05_setup.tex =====
On the quadrotor the
planner updates a commanded position increment once every three $0.01$\,s MuJoCo physics steps.
Between planner updates, the cascaded geometric controller of \autoref{sec:method:geometric}
computes rotor thrusts at every physics step.

% ===== [A1b Table 5.1 X2 row] chapters/05_setup.tex =====
      & spatial position increments
      & $0.03$\,s of simulated time; controller at $0.01$\,s \\

% ===== [A2 obstacle radius] chapters/05_setup.tex =====
the six obstacles of radius $0.025$ and the goal line at $y=0.35$.

% ===== [A3 F01 scoring convention] chapters/05_setup.tex =====
Untightened evaluations are therefore not
reported. An unprojected plan is the same plan on either set, and its
violations are scored against the tightened boundary.
\srcnote{\texttt{config/projection\_eval.yaml}: \texttt{halfspace\_constraints['avoiding-d3il']},
\texttt{enlarge\_constraints: \{'avoiding': 0.025\}}; parity with \texttt{aux\_repo/dpcc} checked in
v2, \autoref{sec:method:env:avoiding}.}

% ===== [A4a aligning constraints vs demos] chapters/05_setup.tex =====
task whose plan is three-dimensional. Neither is respected by the demonstrations, so satisfying them is
left to the projection.

% ===== [A4x box around keep-out (halfspace sanity)] chapters/05_setup.tex =====
the marks alone understate the object --- a
$0.10$\,m box has to be pushed around a $0.06$\,m keep-out region.

% ===== [A4b demonstrations paragraph] chapters/05_setup.tex =====
That segment
crosses the keep-out region in $106$ of the $120$ contexts, and in $119$ of them once the region is
tightened. The demonstrations were recorded without the region
existing, so the task they teach is precisely the one the constraint forbids; going around it is the
projection's work, not something the policy could have imitated.

% ===== [A4c fig:expert-aligning caption] chapters/05_setup.tex =====
red where that straight push crosses the
  circular keep-out region, green where it clears it.

% ===== [A4d fig:expert-aligning dataref] chapters/05_setup.tex =====
\dataref{DA\_in\_Paper/data/expert\_paths.json, extract/expert\_paths.py; contexts from
d3il/environments/dataset/data/aligning/\{train,test\}\_contexts.pkl}

% ===== [A5 s-curve reference route] chapters/05_setup.tex =====
On
UAV-s-curve the
single route satisfies the constraint set the projection later enforces, tightening included: the
generator threads a passage the scene already contains, so nothing in the data has to be corrected. The
projection is then asked to repair a plan that was never broken, which is why that scene can expose a
limit of the tracking controller but cannot order the projection methods.

% ===== [A6 rotor reach scoping (datasets)] chapters/05_setup.tex =====
One thing has to be read correctly in that figure and in every quadrotor result that follows. A violation
is a property of the \emph{path}, not of the aircraft: what is tested is the line of vehicle centres
against the constraint surfaces, and those surfaces already carry the $0.31$\,m rotor reach
(\autoref{tab:eval}), so the vehicle's own extent is inside the boundary rather than measured against it.

% ===== [A7 F03 metric definition] chapters/05_setup.tex =====
    contexts by its median and its minimum, and the share of the initial distance that was closed,
    $(d\sidx{init}-d\sidx{final})/d\sidx{init}$ in per cent, against an initial distance averaging
    $0.4530$\,m: $100\,\%$ is a box brought onto its target and $0\,\%$ one that was not moved towards it
    at all;

% ===== [A8 F08 crossing latched] chapters/05_setup.tex =====
every test-time constraint end; a flight is scored the moment it passes it and flies on as before. On

% ===== [A9 rotor reach scoping (metrics)] chapters/05_setup.tex =====
  \item whether it is violation-free: no control step on which the flown path violates a constraint. The
    test is on the path, not on the airframe: the surfaces it is tested against already carry the
    $0.31$\,m rotor reach (\autoref{tab:eval});

% ===== [A10a multi-seed] chapters/05_setup.tex =====
  \item \textbf{Multi-seed results} --- every D3IL-avoiding table --- are reported as the mean over the

% ===== [A10b single-seed] chapters/05_setup.tex =====
  \item \textbf{Single-seed results} --- the one-seed D3IL-avoiding tables, D3IL-aligning and the
    quadrotor scenes --- carry no seed dimension,

% ===== [A11a tab:train caption] chapters/05_setup.tex =====
The visual
  input is two RGB tensors, each $(N,1,3,96,96)$, encoded to a combined $(N,128)$ feature; $N$ is the
  training batch size.}

% ===== [A11b tab:train header] chapters/05_setup.tex =====
    & & \multicolumn{3}{c}{D3IL-avoiding} & D3IL-aligning & UAV, all scenes \\

% ===== [A11c tab:train prose] chapters/05_setup.tex =====
In the UAV scenes, the state anchor is $(N,9)$ and
the plan is $(N,8,12)$: a spatial position increment, commanded position, measured position and
measured velocity.

% ===== [A12 F12 recipe part of the comparison] chapters/05_setup.tex =====
What varies is how each objective was brought to convergence,
which is a property of the training run and not of the comparison.

% ===== [A13a tab:eval delta row] chapters/05_setup.tex =====
value \ac{DPCC} introduces, and the only quantity the tightened variants change & $0.025$ & $0.025$ & $0.025$ \\

% ===== [A13b tightening prose] chapters/05_setup.tex =====
The \emph{tightening} margin is $0.025$\,m in all three environments, the value \ac{DPCC} introduces. It

% ===== [A14 lookahead order] chapters/05_setup.tex =====
guiding step adds one further evaluation, the lookahead from the projected state, so a plan there costs

% ===== [A15 eta K integral] chapters/05_setup.tex =====
The two families select the same steps except when $\eta\nfe$ is a whole number: the flow-based samplers
project on the last $\lceil\eta\nfe\rceil$ steps, while the diffusion baseline's gate includes the step at
the threshold itself and projects on one more --- eleven rather than ten at $\nfe=20$ and $\eta=0.5$.

% ===== [A16 F05 same plan claim] chapters/05_setup.tex =====
projection runs on $\lceil 1 \rceil = 1$ sampling step and that step is the terminal one --- so
$n\sidx{guide}=0$, endpoint and per-step projection compute the same plan, and a difference between them
there would be noise rather than a property of either method. The same arithmetic gives zero at $\nfe=1$
for every threshold below $1.0$.

% ===== [A17 F05 sampler at every step] chapters/05_setup.tex =====
Where endpoint projection has guiding steps, its sampler is not the plain one on the two average-velocity
models: at every step it queries the network at a zero interval, where the plain sampler queries it at
the step width $1/\nfe$. On those models a difference between endpoint and per-step projection therefore
belongs to the sampler and the point of projection together.

% ===== [B1 guard -> dataref (panels)] chapters/06_results.tex =====
\guard{The four panels added to this figure were rendered

% ===== [B2 frontier band] chapters/06_results.tex =====
value in its panel, so that a fast configuration that reaches the goal less reliably cannot appear as a
favourable trade.

% ===== [B3 smallest guided budget with threshold] chapters/06_results.tex =====
It is read at two budgets. $\nfe=3$ is
the smallest budget at which endpoint projection has a guiding step
(\autoref{sec:res:constraints:degenerate}), and therefore the smallest at which the two methods differ at
all; at $\nfe=10$ both have several.

% ===== [B4 guard ledger path -> dataref] chapters/06_results.tex =====
\guard{No endpoint-projection evaluation of the diffusion baseline exists on D3IL-avoiding, so no
statement in this subsection extends to it. Recorded in
\texttt{data\_status/PENDING\_20260922\_all\_lacking\_runs.md}.}

% ===== [B5 K=3 five of six] chapters/06_results.tex =====
Every $\nfe=3$ point of
\autoref{tab:avoiding-projectors} sits three to eight times to the right of that at the same success, so
each is dominated by a cheaper configuration that solves the task as well. The budget of three is
therefore dropped from this benchmark's operating point, and with it the one budget at which endpoint
projection would be the better projector at a price that matters.

% ===== [B6 bit-identical run named, 2.45x dropped] chapters/06_results.tex =====
Endpoint projection needs sampling steps to guide, and at
one and two evaluations it has none or one (\autoref{sec:method:degenerate},
\autoref{sec:res:constraints:degenerate}); a run at $\nfe=2$ with both methods given the same candidate
plans produced bit-identical rollouts at $2.45$ times the cost.

% ===== [B6b dataref for the bit-identical run] chapters/06_results.tex =====
instantaneous-velocity matching at three evaluations --- and that is the result carried to the tasks that are
run at those budgets, D3IL-aligning and the quadrotor scenes.

% ===== [B5b K=3 recap] chapters/06_results.tex =====
and the budget of three, the first at which
endpoint projection can act, buys no fewer control steps at three to eight times the time, so it is
dominated in the sense of \autoref{sec:res:avoiding:pareto}.

% ===== [B7 F02 avoiding conclusion] chapters/06_results.tex =====
D3IL-avoiding carries the central claim of this thesis, and the results support it: the average-velocity models
Pareto-dominate the diffusion model of \ac{DPCC} on its own benchmark, at its protocol and with its own
temporal U-Net of $4.0$\,M parameters (\autoref{tab:avoiding-conclusion}). CI-MeanFM at one network
evaluation, with temporal consistency and per-step projection, reaches $1.000$ success with constraint
satisfaction as the baseline of record does, in $59.2$ control steps against $70.1$ and at $18.1$\,ms per
step against $553.4$, a thirtieth of the time. Analytic average-velocity matching at the same budget and
projector reaches $0.967$ in $58.6$ steps at $18.3$\,ms, one episode in thirty for the same saving.

% ===== [B8 K=2 guiding steps with threshold] chapters/06_results.tex =====
At one evaluation it has no
guiding step, and at two it has only one; raising the budget

% ===== [B9 F04 constraint set has no bearing] chapters/06_results.tex =====
The generative models are compared on the final box-to-target distance of plans executed
without projection, on which the constraint set has no bearing --- the plan is the same plan --- so all
four models are compared at once; constraint violations and the time one action costs are compared in
\autoref{sec:res:aligning:projection}.

% ===== [B10 moves the box towards] chapters/06_results.tex =====
One model moves the box to its target on this task and the other three barely move it at all

% ===== [B11a Table 6.5 caption] chapters/06_results.tex =====
contexts from an initial $0.4530$\,m, in brackets the share of that distance closed
  (\autoref{sec:setup:metrics:aligning});

% ===== [B11b reading rule] chapters/06_results.tex =====
The outcome axis is the share of the starting distance
the box was moved towards its target, the bracketed number of \autoref{tab:va-models}, so the dashed
line at $0\,\%$ is a box that was not moved and $100\,\%$ at the top is a box delivered;

% ===== [B11c fig:aligning-tradeoff caption] chapters/06_results.tex =====
\caption[Alignment: outcome against cost]{D3IL-aligning without projection: the share of the
  starting distance the box was moved towards its target, median over the ten contexts, against the time
  to compute one action on a log axis; every cell of \autoref{tab:va-models} except the two
  $\alpha_{\mathrm{end}}=0.05$ rows, one training seed. The
  dashed red line at $0\,\%$ is a box not moved. Rings mark the non-dominated configurations, joined by
  the dashed staircase; beside each point is its step budget; the boxed arrow points where both axes
  improve. The outcome axis is logarithmic in the distance left to the
  target; the tick labels are the share closed.}

% ===== [B11d Table 6.7 caption] chapters/06_results.tex =====
contexts, the median with the share of the initial $0.4530$\,m closed in brackets, the mean $\pm$

% ===== [B11e section 6.4 percent] chapters/06_results.tex =====
The median context ends $0.197$\,m from the target, $57\,\%$ of the initial distance closed.

% ===== [B11f fig:aligning-projected-tradeoff caption] chapters/06_results.tex =====
the faint dashed grey ring is the
  unprojected MeanFM cell at the same budget. Filled points keep at least nine contexts violation-free,
  hollow ones do not; rings and the dashed staircase mark the non-dominated filled points; the dashed line at $0\,\%$ is a box not moved.}

% ===== [B11g comment on the label] chapters/06_results.tex =====
% dashed line: under it, the bottom frame struck it through (already so before v3.77).

% ===== [B12 context numbering in the guard] chapters/06_results.tex =====
obstacle until it overlaps the box of context~8, the evaluation's box--obstacle guard holds the box, and

% ===== [B13 editorial sentence out of the guard] chapters/06_results.tex =====
that one. Comparisons between rows are unaffected. Whether the chapter states it is the author's call
(\texttt{DA\_20260923\_R2fix\_R37\_aligning.md} \S6.2).}

% ===== [B14 F04 twin ring] chapters/06_results.tex =====
Every point of the second is a plan \emph{after} projection on the tightened
set, so it is dearer per control step and, for the same model and budget, further from the target than its
unprojected twin --- which is drawn as a faint grey ring for analytic average-velocity matching at the same
budget, so the price can be read directly: from $84\,\%$ of the distance closed at $190$\,ms to
$61\,\%$ at $266$\,ms.

% ===== [B14b which constraint is crossed (halfspace sanity)] chapters/06_results.tex =====
the unprojected plan ending closest of all while crossing the constraint in eight
contexts of ten.

% ===== [B15 pillars same plans] chapters/06_results.tex =====
execute the
same plans instead of the manipulator:

% ===== [B16 changes nothing] chapters/06_results.tex =====
Before projection the vehicle changes nothing (\autoref{tab:uav-pillars-raw}, upper block).

% ===== [B17 guarantees] chapters/06_results.tex =====
so what the per-step projector guarantees on D3IL-avoiding it guarantees on UAV-pillars. Yet success falls

% ===== [B18 tracking percentile] chapters/06_results.tex =====
The vehicle tracks more
closely: its flown position stays within about $0.007$ of the commanded one in the table's units, against the

% ===== [B19 corridor band dominance] chapters/06_results.tex =====
Among the configurations with the fewest
violating steps, instantaneous-velocity matching at twenty evaluations is on the frontier and the baseline is
not: under the tilt per-step projection, $252.6$ steps at $360$\,ms, dominates endpoint projection ($269.8$,
$417$\,ms) and the baseline ($290.3$, $666$\,ms); over the hump endpoint projection with one candidate is the
cheaper, $316.2$ steps at $291$\,ms, per-step projection the shorter, $285.1$ at $355$\,ms, and the baseline,
$320.8$ steps at $654$\,ms, is dominated.

% ===== [B20 pillars same plan flown] chapters/06_results.tex =====
its result is the difference between a plan's outcome on
D3IL-avoiding and the outcome of the same plan flown,

% ===== [B21 corridor within 1.6] chapters/06_results.tex =====
with it they leave the same violating steps at a shared
budget and projection method, to within $1.6$ per flight, while under the tilt the budget takes them from

% ===== [B22 F08 flight endings] chapters/06_results.tex =====
A flight ends in one of
three ways: it crosses the finish line, it reaches the episode limit, or the divergence guard ends it
because the vehicle inverted (\autoref{sec:setup:metrics:uav}).

% ===== [B23 Table 6.13 caption success] chapters/06_results.tex =====
training seed. \emph{Success}: crosses the finish line; \emph{violating steps}: per flight, counted until

% ===== [B24 keeping the plan] chapters/06_results.tex =====
What the controller changes is tested by
keeping the plan and changing the controller:

% ===== [B25 Table 6.15 title/caption] chapters/06_results.tex =====
\caption[UAV-s-curve: the same plans under two tracking controllers]{UAV-s-curve, FM at $\nfe=1$, the
  configuration selected in \autoref{tab:uav-scurve}, the same plans under the two tracking controllers,

% ===== [B26 stale pilot sentence] chapters/06_results.tex =====
The
controller that flies this scene is about twenty-four times the price of the one that does not, and by itself
it costs twelve times as much per control step as generating the plan

% ===== [B26b isolates the controller] chapters/06_results.tex =====
The remainder row therefore carries the
simulator step with the controller, and it is only their \emph{difference} that isolates the controller.
One training seed, ten flights against ten.}

% ===== [B27b same projected plans] chapters/06_results.tex =====
and the same projected plans succeed on five flights of ten or on none, and
the same unprojected plans violate on $23$ control steps per flight or on $42$, depending on which controller
flies them.

% ===== [B27a twenty-fourth (caveat)] chapters/06_results.tex =====
fewer violating
steps in both, at a twenty-fourth of the cost per control step. MuJoCo MPC's one gain, keeping the unprojected

% ===== [B28a same plans (conclusion)] chapters/06_results.tex =====
on a route of two sharp turns the same plans succeed or fail by the controller that
flies them.

% ===== [B27c twenty-fourth (conclusion)] chapters/06_results.tex =====
succeeds on none, at a twenty-fourth of MuJoCo MPC's cost per control step. MuJoCo MPC's one gain, keeping

% ===== [B29a follows from the data] chapters/06_results.tex =====
That the
simplest objective leads follows from the data:

% ===== [B28b guard same plans] chapters/06_results.tex =====
not on
the controller comparison, which flies the same plans under both controllers.}

% ===== [B30a 6.4 intro] chapters/06_results.tex =====
On the two D3IL tasks the flow-based models
Pareto-dominate the diffusion model of \ac{DPCC} with its own temporal U-Net of $4.0$\,M parameters; on the
quadrotor they carry that result to a new plant,

% ===== [B30b 6.4 avoiding item] chapters/06_results.tex =====
The average-velocity models Pareto-dominate the baseline.
    CI-MeanFM at one evaluation holds

% ===== [B28c 6.4 s-curve item same plans] chapters/06_results.tex =====
The same plans succeed or fail by the controller that flies
    them, and the cascaded geometric controller is the better one:

% ===== [B27d 6.4 s-curve twenty-fourth] chapters/06_results.tex =====
projection, five of ten against none with it, at a twenty-fourth of the cost. Instantaneous-velocity

% ===== [B30c 6.4 conclusion three steps] chapters/06_results.tex =====
there the
average-velocity models Pareto-dominate the diffusion model of \ac{DPCC} with its own temporal U-Net, and no

% ===== [B30d 6.4 conclusion picture] chapters/06_results.tex =====
the average-velocity models lead: on D3IL-avoiding they Pareto-dominate the diffusion model of
\ac{DPCC} with its own temporal U-Net, in fewer control steps at a thirtieth of its time per action, and on

% ===== [B29b as expected] chapters/06_results.tex =====
on UAV-corridor the three
objectives cannot be told apart, as expected from what the average-velocity objectives are built for, and on

% ===== [B31 any flow-based objective] chapters/06_results.tex =====
  \item[UAV-corridor.] Any flow-based objective, the budget chosen for the trade: violating steps per flight

